% Layout advisor: predicting the best layout, page size and execution options from a calibrated
% cost model, validated on a second machine. Figures and tables: scripts/report_advisor.py;
% data: scripts/run_report.sh (results/report/, results/ground_truth/laptop_v1).
\section{Choosing the Layout Without Running the Queries}
\label{sec:advisor}

Sections~4--7 explain each result after measuring it. Here we ask the reverse question:
\emph{given a machine, a table and a weighted workload, can a model choose the layout, the PAX
page size and each query's execution options without running the queries?} Our advisor
calibrates the machine once, describes every query on every layout as the memory accesses its
loop makes, predicts the time of each combination and searches the candidate space. We
evaluate it against 10,368 measurements on a second machine.

\subsection{Engine and machine}
The machine is an Intel Core i7-1360P (Raptor Lake) under native Linux~6.8. We use one
performance core, with 48\,KB 12-way L1D, 1.25\,MB 10-way L2, 18\,MB 12-way shared L3 and
64-byte lines, built with \code{-O2 -march=native}. The engine is an earlier branch of ours.
It has the same data model and layouts as Section~2, and runs queries as follows:
\begin{itemize}
  \item \textbf{NSM} works tuple at a time. Predicates either short-circuit (one branch per
  predicate) or are evaluated branch-free, and contiguous projections are copied with
  \code{memcpy}.
  \item \textbf{DSM} processes vectors of $B$ rows ($B$ from 256 to 128K). The first predicate
  builds a selection vector, later predicates refine it at the surviving positions, and the
  projected columns are then gathered: Section~2's \emph{late} strategy.
  \item \textbf{PAX} does the same one page at a time.
\end{itemize}
Output tuples go into a 1024-tuple buffer that is folded into a checksum, so writes stay in
cache. Reconstruction therefore costs gathering, not write-allocate. The options the advisor
chooses per query are $B$, branchy or branch-free predicates, and predicate order.

The layout results of Section~4 carry over to this machine, though the gaps are smaller:
\begin{itemize}
  \item From a 1\,GB table, DSM's \code{SUM(c0)} is 12.6$\times$ faster than NSM's (15$\times$
  on the Ryzen).
  \item PAX matches DSM only with pages of 4\,MB or more (0.47\,ns per row for both); 16\,KB
  pages are 3.1$\times$ slower.
  \item For the sum of all 16 attributes from DRAM, NSM beats DSM from 20\% selectivity on
  (1.03--1.13$\times$).
\end{itemize}

\begin{table}[t]
\caption{The second machine as calibrated by \code{build/calibrate} (one P-core; bandwidth
of a sequential read, latency of a random pointer chase).}
\label{tab:advhw}
\small\gen{advisor_hw}
\end{table}

\subsection{Calibration}
Sysfs reports cache sizes, but not what a miss costs, so \code{build/calibrate} measures every
constant on the pinned core (\cref{tab:advhw}). Each sample lasts at least 20\,ms, and we take
the median of 5. The measurements are:
\begin{itemize}
  \item \textbf{Levels.} Sequential reads and a random pointer chase over regions from 4\,KB
  to 1\,GB give each level's bandwidth, latency and capacity knee.
  \item \textbf{Transfer unit.} Strided loads give $B_\mathrm{eff}$, the DRAM bytes moved per
  line miss. It is 156\,B, not 64, because the adjacent-line prefetcher fetches pairs of
  lines.
  \item \textbf{MLP.} Independent pointer chains give the number of misses that overlap.
  \item \textbf{TLB.} Chases on 4\,KB and 2\,MB pages give the TLB reach and the page-walk
  cost.
  \item \textbf{Prefetcher.} Runs separated by gaps give a table of how many bytes the
  prefetcher over-fetches.
  \item \textbf{CPU.} The engine's own kernels, run on in-cache data, give the cycles per row,
  per gathered value, per vector and per page, and the cost of a mispredicted branch.
\end{itemize}
The whole calibration takes 35\,s, and no machine constant appears in the code.

\subsection{Cost model}
Following Manegold et al.~\cite{manegold}, the advisor describes a query on a layout as a
composition of basic access patterns:
\begin{itemize}
  \item a sequential pass over a region;
  \item $n$ independent random accesses, used for point lookups;
  \item a small buffer that is reused, such as the selection vector or the output buffer.
\end{itemize}
Patterns are combined sequentially, or concurrently, in which case they share each cache in
proportion to their footprints. One function per layout produces the pattern from the
engine's loop.

The central term generalises Section~4's $f(s)$. Let $f$ be the fraction of rows a column is
read at, and $G$ the mean run of skipped rows. A line holding $v$ candidate values is then
touched with probability
\[ p = 1-(1-f)\,(1-1/G)^{\,v-1}. \]
For independent rows, $G = 1/f$ and this reduces to $1-(1-s)^v$. On sorted or clustered data,
qualifying rows come in long runs and far fewer lines are touched. $G$ comes from a data
profile (histograms, sortedness, run lengths) that also estimates each predicate's selectivity.

A pass misses once on every line it touches, because each query starts cold and never comes
back to a line. A touched line counts as sequential if its neighbour is touched too, and as
random otherwise. Memory time then sums, over the levels $i$,
\[ T_\mathrm{mem}=\sum_i M^s_i\,B\Big(\frac{1}{W_{i+1}}-\frac{1}{W_i}\Big)
 + M^r_i\,\frac{\lambda_{i+1}-\lambda_i}{\mathrm{MLP}}, \]
where $M^s_i$ and $M^r_i$ are the sequential and random misses at level $i$, $W$ is bandwidth
and $\lambda$ is latency. TLBs are modelled as caches with 4\,KB lines.

Four corrections bring the model from the published one to a modern machine:
\begin{itemize}
  \item misses at the last cache level are counted in units of $B_\mathrm{eff}$;
  \item random misses overlap, so each costs latency divided by the MLP;
  \item short runs pay the calibrated prefetch over-fetch;
  \item page walks inside sequential streams are hidden.
\end{itemize}
CPU time $C$ is calibrated cycles times counts: predicate evaluations per row and per
survivor, gathered values and copied tuples. To it we add the branch misses $\mathit{Br}$
($2r(1-r)$ per evaluated row, or 0 when the data is sorted) and the per-vector and per-page
overheads $O$. Finally
\[ T=\max(k_m T_\mathrm{mem}, T_c)+\gamma\min(k_m T_\mathrm{mem}, T_c),\qquad
   T_c=k_c C+k_b \mathit{Br}+k_o O . \]
Each of the six parameters ($k_m$, $k_c$, $k_b$, $k_o$, $\gamma$, and the cost of locating a
lookup's PAX page) scales one physical term. They are fitted by least squares on
$\log(\text{predicted}/\text{measured})$.

\textbf{Search.} For each of 12 storage configurations (NSM, DSM, and PAX with pages from
1\,KB to 16\,MB), every query gets its cheapest options. This per-query choice is exact, because
the score is a weighted sum and every query starts cold. The configurations are then ranked by
the weighted mean time. The advisor reports:
\begin{itemize}
  \item the runner-up, and a close call when the margin is under 5\%;
  \item whether the choice still holds with every selectivity scaled by 0.25--4$\times$ and
  memory 20\% slower or faster.
\end{itemize}
A search makes 72 predictions in 20\,ms.

\subsection{Methodology}
\textbf{Data and workloads.} The ground truth uses three 256\,MB tables, each 14$\times$ the
L3: 16 attributes, 64 attributes, and 16 attributes with a sorted \code{c0}. \code{c0} has
1000 distinct values, so selectivities down to 0.1\% are exact. There are four query families:
\begin{itemize}
  \item scans;
  \item one predicate at 10 selectivities from 0.1\% to 100\%;
  \item a second predicate at 50\%, in both evaluation orders;
  \item 1K--100K point lookups, uniform or sorted.
\end{itemize}
Each family uses four projections: \code{c0}, the first 25\% of the attributes, all of them,
and every fourth one. Each query is measured on all 12 storage configurations with every
applicable option, 10,368 points in all.

\textbf{Protocol.} Every point is measured cold: the table is flushed with \code{clflush}
before each sample. Each point gets 1 warm-up and 5 timed samples, and we report the median.
Results are checked against a reference first, as in Section~2. The run takes 100\,min. A
repeat of 10\% of the points differs by a median of 4.3\% (p90 19.6\%), and 30\% of queries
have a runner-up within 5\%. This noise bounds any decision metric.

\textbf{Splits and metrics.}
\begin{itemize}
  \item We fit on 30\% of the queries of the 16-attribute uniform table (1,053 points) and
  report errors on the other 9,315, which include both other tables.
  \item Decisions are tested on 1000 random workloads. Each has 2--6 queries from one table,
  with weights drawn log-uniformly from 1 to 1000.
  \item A choice is top-1 if its measured time is within 5\% of the best. Regret is the
  measured time of the choice divided by the best measured time.
  \item The advisor is charged at the options it picks itself, while fixed rules are given the
  best measured options.
\end{itemize}

\begin{table}[t]
\caption{Model ladder on the 9,315 held-out points (error = predicted $\div$ measured $- 1$)
and the decisions each rung makes on 1000 random workloads (regret).}
\label{tab:ladder}
\small\gen{advisor_ladder}
\end{table}

\begin{figure*}[t]
\centering
\fig[0.92\linewidth]{advisor_accuracy}
\caption{(a) Predicted against measured latency of the held-out points, by layout.
(b) Distribution of regret over 1000 random workloads for the advisor and four fixed rules
(higher and further left is better).}
\label{fig:advisor}
\end{figure*}

\subsection{Results}
\textbf{Accuracy.} \Cref{tab:ladder} adds one group of mechanisms per rung:
\begin{itemize}
  \item \textbf{v1}, the model as published, has a p90 error of 618\%. Every random line costs
  a full DRAM latency, and every miss moves 64 bytes, although the hardware overlaps about 15
  misses and moves up to 156 bytes per line.
  \item \textbf{v2} adds $B_\mathrm{eff}$, MLP, over-fetch and hidden walks, which cut the
  median error to 25\%.
  \item \textbf{v3} adds branches, per-vector and per-page overheads and overlap, giving 16\%
  (p90 45\%; \cref{fig:advisor}a).
\end{itemize}
Predicted DRAM bytes match the memory-controller counters: predicted $\div$ measured is 0.98
for NSM, 0.89--0.92 for DSM and 0.87 for PAX.

\begin{table}[t]
\caption{Decisions on 1000 random workloads: the advisor against fixed rules (which are given
the best measured options).}
\label{tab:decisions}
\small\gen{advisor_decisions}
\end{table}

\textbf{Decisions.} Always-DSM is the best fixed rule, and also the best single choice in
hindsight for every table. It is right more often than the advisor (86\% against 72\%,
\cref{tab:decisions}), but it is 8.35$\times$ slower than the best on lookup-heavy mixes. The
advisor's worst case is 1.51$\times$, and its median regret is 1.08 (\cref{fig:advisor}b). Judged
on its storage choice alone, its median regret is 1.00 and its maximum 1.49.

\textbf{Generalisation and robustness.}
\begin{itemize}
  \item Fitted on the 16-attribute uniform table without lookups, the model's median error is
  20\% on the unseen 64-attribute table and 16\% on the unseen sorted one. The decisions do
  not change (top-1 72\%, worst case 1.54$\times$).
  \item Scaling any group of calibrated constants by $\pm$20\%, or the MLP by 0.5--2$\times$,
  changes 0.5--7\% of the decisions.
  \item Scaling the selectivity estimates by 0.5 or 2 changes 13--15\%, so the estimates
  matter more than the hardware constants.
\end{itemize}

\textbf{Example.} On a 4\,GB table of 8M rows $\times$ 64 attributes, the advisor picks:
\begin{itemize}
  \item DSM for narrow aggregates, with PAX at 4\,MB pages 1.03$\times$ behind;
  \item PAX with 4\,MB pages for a mix that is half 1\% \code{SELECT *} and half aggregates
  and lookups. DSM is 1.01$\times$ behind and NSM 2.07$\times$, and the advisor flags the
  choice as a close call that holds only for selectivity scaled by 0.8--2$\times$.
\end{itemize}

\subsection{Analysis}
\textbf{What the model gets right.} The line count of Section~4 decides the ranking, but its
cost per line depends on the access pattern:
\begin{itemize}
  \item A whole NSM row is 128\,B, and the adjacent-line prefetcher fetches it whole. The model
  therefore predicts 129\,B per row from DRAM for any $k$, against 105--148\,B measured.
  Counting 64-byte lines underestimates this by up to 2$\times$.
  \item A DSM gather touches $f(s)$ of each column's lines, but at random. Those lines are
  cheap only because about 15 misses overlap, so without MLP the model would make late
  materialization look many times slower than it is.
\end{itemize}

\textbf{Where it fails.}
\begin{enumerate}
  \item \textbf{Small PAX pages on the 64-attribute table} are over-predicted by up to
  8$\times$. Footprint-proportional cache sharing gives the 8-byte per-page selection vector a
  3-byte share of L1, which produces phantom misses. In addition, $k_o = 15.6$ inflates the
  per-page overhead.
  \item \textbf{Unfiltered reconstruction} on DSM and PAX is over-predicted by 34--45\%.
  \item \textbf{Point lookups} are under-predicted by 19--28\%. This appeared once full-size
  calibration corrected the page-walk cost to 11\,ns; the remaining cost is not yet modelled.
\end{enumerate}

\textbf{Wrong picks.} Most of the advisor's wrong picks cost little. 157 of its 283 wrong
choices are PAX with 1\,MB pages where DSM was best, and these two layouts are within the
noise of each other for most queries. Top-1 accuracy is therefore limited by the 4--20\%
measurement noise and the many near-ties. The advisor's real value is that it avoids the
4--12$\times$ losses every fixed rule suffers on some workloads.

\textbf{Summary.} A model built from cache lines and calibrated on a single core in 35\,s
predicts latency to within 16\% (median). Its picks are on average within 8\% of the best, and
never more than 1.51$\times$ worse, across layouts, page sizes and execution options. Both the
model and the search run on another machine after recalibration (\code{make};
\code{scripts/run\_report.sh}). We have measured only on this one; checking that is future
work.
